% No \label here. \exercises is a starred section, so a label attached to it
% takes the number of the last *numbered* section instead (quyet dinh 68).
% Prose refers to these by name.
\exercises

\tierunderstand

\begin{bt}
  \item Dẫn lại bảng ba chế độ của mục~\ref{sec:ch10-hai-rang-buoc} cho một
    tầng tổng quát: ảnh \(H \times W\) với \(C_{\text{in}}\) kênh vào,
    \(C_{\text{out}}\) kênh ra, kernel \(K \times K\), giữ nguyên độ phân giải.
    Viết ba công thức số tham số. Rồi trả lời câu đắt hơn: vì sao phép chia do
    \tn{chia sẻ trọng số}{weight sharing} đúng bằng \(HW\) mà không phụ thuộc
    \(K\), trong khi phép chia do \tn{tính cục bộ}{locality} thì phụ thuộc cả
    hai?

  \item Mục~\ref{sec:ch10-hai-rang-buoc} đo được 71.91\% tham số tự do của tầng
    H1 trong bài 1989 là bias. Tính lại con số ấy cho một tầng
    \code{nn.Conv2d} hiện đại có bias dùng chung theo kênh. Rồi nói: nếu bài
    1989 dùng bias dùng chung thì tổng 9,760 thành bao nhiêu, và tỷ số 6.625
    thành bao nhiêu?

  \item LeNet-5 dựng lại ra đúng 340,908 kết nối chỉ khi đếm cả tầng ra RBF, vốn
    có 840 kết nối và không tham số tự do nào. Nêu một cách đọc khác của cụm
    \enquote{trainable free parameters} làm cho con số 60,000 vẫn đúng mà cách
    chia tầng thì khác, rồi nói phép kiểm nào tách được hai cách đọc bằng đúng
    những con số bài in.

  \item Chứng minh \eqref{eq:ch10-dang-bien} cho \tn{tích chập}{convolution} một
    chiều với padding vòng, thẳng từ định nghĩa của phép ấy. Rồi chỉ ra bước nào
    trong chứng minh hỏng khi padding là padding không, và bước nào hỏng khi có
    stride bằng 2.

  \item Mục~\ref{sec:ch10-dang-bien} đo được một mạng ba tầng pooling
    \tn{đẳng biến}{equivariance} ở bội của 8. Tổng quát hóa: cho một mạng có các
    tầng stride
    \(s_1, \dots, s_L\), chu kỳ đẳng biến là bao nhiêu? Rồi tính cho ResNet-50,
    tra stride từng tầng của nó, và nói con số ấy so với 224 pixel đầu vào thì
    là bao nhiêu phần.

  \item Định lý 1 của Cordonnier và cộng sự đòi \(N_h = K^2\) head cho kernel
    \(K \times K\). Mục~\ref{sec:ch10-attention-lam-duoc} nói self-attention vì
    thế \enquote{biểu diễn được} tích chập. Viết ra ba thứ mà phát biểu ấy
    \emph{không} nói, mỗi thứ kèm một phép đo phân biệt được nó với phát biểu đã
    chứng minh.

  \item Mục~\ref{sec:ch10-anh-that} đo được khoảng cách CNN với MLP rộng ra theo
    lượng dữ liệu, còn bội số ảnh cần thì cho kết luận ngược. Dựng một ví dụ số
    nhỏ - hai đường cong độ chính xác theo lượng dữ liệu, viết bằng công thức -
    trong đó khoảng cách tuyệt đối rộng ra đơn điệu \emph{và} bội số ảnh cần
    hẹp lại đơn điệu. Rồi nói điều đó nghĩa là gì cho việc đọc bất kỳ bảng
    ablation nào.

  \item Ở phép dịch 8, mục~\ref{sec:ch10-anh-that} đo CNN đổi nhãn trên 59.81\%
    số ảnh, và nêu ba nguyên nhân rời nhau: mạng pad bằng không, bộ phân loại
    đọc vị trí, và phép dịch vòng thật sự đổi ảnh. Thiết kế một phép đo tách
    được ba phần ấy. Nói rõ mỗi phần cần đối chứng nào, và phần nào trong ba
    phần bạn khử được mà không phải huấn luyện lại mô hình.
\end{bt}

\tierapply

Trạng thái xuất phát cho cả năm bài: clone repo companion, đứng ở thư mục gốc
của nó, và checkout đúng tag của chương này.

\begin{minted}{console}
$ git clone https://github.com/Giang-Dang/rnn-to-transformer-lab.git
$ cd rnn-to-transformer-lab
$ git checkout ch10
$ conda env create -f environment.yml
$ conda activate rnn-to-transformer-lab
$ python -c "from rnn_to_transformer_lab.cifar import fetch; fetch()"
$ python verify.py --only ch10
\end{minted}

Dòng áp chót là bắt buộc và chỉ chạy một lần: nó tải CIFAR-10 về rồi cache
lại. Trên máy tôi việc ấy mất ba mươi hai phút, nên đừng tưởng là treo. Nó
không in gì trong lúc tải; muốn biết nó còn sống thì xem kích thước file trong
thư mục \code{data/} lớn dần lên. Mọi lần chạy sau đọc cache và không cần mạng.

Dòng cuối phải in ra \code{verify: ok}. Bốn mục của chương này chạy hết chừng
năm phút cộng lại, và gần hết chỗ ấy là một mục: phép quét ở
mục~\ref{sec:ch10-anh-that} huấn luyện 36 mô hình. Ba mục còn lại không huấn
luyện gì và xong trong vài giây.

Vẫn cảnh báo cũ: đừng gọi các script qua \code{conda run}.

\begin{bt}
  \item \textbf{Đo cái mục~\ref{sec:ch10-con-lai-gi} nói là chưa ai đo.} Bảng
    của chương so CNN với MLP, tức so hai ràng buộc cùng lúc. Dựng mô hình ở
    giữa: cục bộ mà \emph{không} dùng chung trọng số, tức mỗi vị trí có kernel
    riêng của nó. PyTorch không có sẵn tầng ấy; dựng bằng \code{unfold} cộng một
    phép nhân theo lô, hoặc bằng cách khác bạn thấy gọn hơn. Rồi chạy lại đúng
    phép quét của \code{ch10\_cifar.py} với ba mô hình thành bốn, và trả lời:
    trong bội số 28.4 lần, bao nhiêu thuộc về tính cục bộ và bao nhiêu thuộc về
    chia sẻ trọng số? Cảnh báo trước: mô hình ấy có rất nhiều tham số, nên phép
    so \enquote{cùng số tham số} sẽ không dựng được cùng lúc với phép so
    \enquote{cùng cấu trúc}, và bạn phải chọn một rồi nói mình chọn cái nào.

  \item Sửa \code{equivariance\_error} cho nó nhận stride không phải lũy thừa
    của 2, rồi quét: với tầng stride \(s\), tập các phép dịch đầu vào cho sai
    số đúng bằng không là tập nào? Đoán trước rồi chạy. Rồi chạy bộ test của
    chương:

    \begin{minted}{console}
$ python -m pytest -q tests/test_ch10.py
    \end{minted}

    Test nào bắt được thay đổi của bạn, test nào không, và vì sao?

  \item Đẩy phép dựng của mục~\ref{sec:ch10-attention-lam-duoc} về phía thật.
    Thay các head đặt bằng tay bằng một tầng self-attention có mã hóa vị trí
    tương đối bậc hai như bài Cordonnier mô tả, rồi tăng dần hệ số làm softmax
    nhọn lên và đo sai số so với \code{F.conv2d} theo hệ số ấy. Vẽ đường cong.
    Nói rõ đường cong ấy đo cái gì, và vì sao nó \emph{không} phải bằng chứng
    rằng một mô hình được huấn luyện sẽ đi tới đó.

  \item Bảng ở mục~\ref{sec:ch10-anh-that} chuẩn hóa dữ liệu bằng thống kê của
    chính tập huấn luyện chứ không của cả 50,000 ảnh. Đổi sang dùng thống kê
    toàn bộ, chạy lại, và đo xem hàng 1000 ảnh đổi bao nhiêu. Rồi trả lời: chỗ
    rò rỉ ấy làm bảng đẹp lên hay xấu đi, và vì sao nó lại tác động mạnh nhất
    lên đúng hàng mà chương đọc kỹ nhất?

  \item Dựng lại bảng số tham số của chương cho một kiến trúc hiện đại bạn tra
    được đủ siêu tham số - ResNet-18 là lựa chọn rẻ. Tính tỷ lệ tham số nằm ở
    các tầng chia sẻ so với các tầng kết nối đầy đủ, rồi đặt cạnh 0.6250 của bài
    1989 và 0.9714 của LeNet-5. Con số đi theo hướng nào, và cái gì trong lịch
    sử kiến trúc giải thích hướng ấy?
\end{bt}

\tierextend

\begin{bt}
  \item Mục~\ref{sec:ch10-anh-that} đo bội số 28.4 lần trên CIFAR-10 với một
    mạng ba tầng tích chập. Con số ấy là của cặp mô hình ấy trên bộ dữ liệu ấy,
    và chương không tổng quát hóa nó. Câu hỏi: bội số ấy phụ thuộc cái gì? Thiết
    kế một phép đo cho \emph{một} biến - độ phân giải ảnh, số lớp, mức độ ảnh
    thật sự có cấu trúc cục bộ - và chạy nó ở quy mô bạn chạy được. Nói trước
    kết quả nào sẽ làm bạn nghi ngờ chính thiết kế của mình. Gợi ý cho biến thứ
    ba: xáo cố định thứ tự pixel của mọi ảnh bằng cùng một hoán vị phá cấu trúc
    cục bộ mà không đổi thông tin, và MLP không được ảnh hưởng gì.

  \item Bài Cordonnier chứng minh sức biểu diễn và đo được rằng các tầng đầu của
    mạng chỉ-attention tự học ra kiểu chú ý theo lưới. Chưa ai trả lời câu ở
    giữa: \emph{bao nhiêu} dữ liệu thì đủ để nó học ra? Dựng phép đo ấy. Huấn
    luyện một mạng self-attention nhỏ trên ảnh ở vài lượng dữ liệu, đo mức độ
    phân phối attention của tầng đầu giống một kernel cục bộ - bạn phải tự định
    nghĩa độ đo ấy và bảo vệ định nghĩa - rồi đặt đường cong ấy cạnh đường cong
    độ chính xác. Hai đường có gãy ở cùng chỗ không?

  \item Mục~\ref{sec:ch10-dang-bien} đo được phép \tn{giảm mẫu}{subsampling} phá
    đẳng biến ở đúng một nửa số phép dịch, và Zhang đề xuất lọc chống aliasing
    để vá. Đọc cách vá
    ấy như một phát biểu về đánh đổi chứ không như một kỹ thuật: nó mua lại tính
    chất gì, và trả bằng gì? Cài nó, đo cả hai vế trên bảng của
    mục~\ref{sec:ch10-anh-that}, và trả lời câu mà bài không hỏi: nếu đẳng biến
    ở mọi phép dịch đáng giá đến thế, vì sao kiến trúc thắng giải năm nào cũng
    vẫn có stride?

  \item Chương này và chương~\ref{ch:tien-huan-luyen} đo hai thứ mà cả hai đều
    gọi là \enquote{đói dữ liệu}, ở hai quy mô cách nhau năm bậc độ lớn. Đề xuất
    một đại lượng duy nhất đo được ở cả hai chỗ, định nghĩa nó đủ chặt để tính
    được từ những bảng hai chương đã in, rồi tính. Nếu bạn kết luận là không có
    đại lượng nào như vậy, hãy chứng minh chỗ nào cản - và đó cũng là một câu
    trả lời, miễn là bạn nói được cản ở đâu.
\end{bt}
